\begin{anexosenv}

% Imprime uma página indicando o início dos anexos
\partanexos

% ===========================================================================
\chapter{Categorias de Gravidade de Dano NCC MERP}
\label{anex_ncc_merp}
% ===========================================================================

O \autoref{qua_ncc_merp_completo} apresenta a taxonomia completa das
9 Categorias de Gravidade de Dano da \textit{National Coordinating Council
for Medication Error Reporting and Prevention} (NCC MERP) \cite{nccmerp2001},
conforme implementadas no catálogo de referência do SIGEA-GTT.

\begin{quadro}[htb]
\centering
\small
\caption{Categorias de Gravidade de Dano NCC MERP (A a I).}
\label{qua_ncc_merp_completo}
\begin{tabularx}{\textwidth}{c c X c}
\toprule
\textbf{Categoria} & \textbf{Nível} & \textbf{Descrição} & \textbf{Dano Real} \\
\midrule
\textbf{A} & Potencial &
  Circunstâncias ou incidentes com capacidade de causar erro,
  mas que não atingem o paciente. & Não \\
\midrule
\textbf{B} & Nenhum &
  Ocorreu um erro, mas não atingiu o paciente. & Não \\
\midrule
\textbf{C} & Nenhum &
  Ocorreu um erro que atingiu o paciente, mas não lhe causou dano. & Não \\
\midrule
\textbf{D} & Nenhum &
  Ocorreu um erro que atingiu o paciente e necessitou de monitoramento
  adicional para confirmar ausência de dano, ou intervenção para
  prevenção do dano. & Não \\
\midrule
\textbf{E} & \textbf{Dano Temporário} &
  Ocorreu um erro que contribuiu para ou resultou em dano temporário
  ao paciente e necessitou de intervenção. & \textbf{Sim} \\
\midrule
\textbf{F} & \textbf{Dano Temporário} &
  Ocorreu um erro que contribuiu para ou resultou em dano temporário
  ao paciente e necessitou de hospitalização inicial ou prolongada. & \textbf{Sim} \\
\midrule
\textbf{G} & \textbf{Dano Permanente} &
  Ocorreu um erro que contribuiu para ou resultou em dano permanente
  ao paciente. & \textbf{Sim} \\
\midrule
\textbf{H} & \textbf{Quase Fatal} &
  Ocorreu um erro que necessitou de intervenção necessária para a
  manutenção da vida do paciente (ex: ressuscitação cardiopulmonar). & \textbf{Sim} \\
\midrule
\textbf{I} & \textbf{Fatal} &
  Ocorreu um erro que pode ter contribuído para ou resultou na morte
  do paciente. & \textbf{Sim} \\
\bottomrule
\end{tabularx}
\fonte{Adaptado de \citeonline{nccmerp2001}.}
\end{quadro}

% ===========================================================================
\chapter{Resolução CNS 466/2012 -- Ética em Pesquisa com Dados Clínicos}
\label{anex_cns_466}
% ===========================================================================

O desenvolvimento do SIGEA-GTT observou rigorosamente os preceitos éticos
estabelecidos pela Resolução do Conselho Nacional de Saúde (CNS)
n.º 466, de 12 de dezembro de 2012 \cite{cns466_2012}, que regulamenta
pesquisas envolvendo seres humanos no Brasil, e pela Lei Geral de Proteção
de Dados (LGPD -- Lei n.º 13.709/2018) \cite{lgpd2018}.

Os principais cuidados éticos adotados no projeto foram:

\begin{enumerate}
    \item \textbf{Não utilização de dados reais de pacientes}: Todos os
    prontuários presentes no sistema são \textbf{fictícios}, gerados
    exclusivamente para fins didáticos e de treinamento, sem nenhuma
    vinculação com pacientes reais do Hospital Universitário da UFS ou
    qualquer outra instituição de saúde;

    \item \textbf{Isolamento completo de bases}: O SIGEA-GTT não estabelece,
    em nenhuma hipótese, conexões com bases de dados de produção do HU-UFS,
    do EBSERH, do DATASUS ou de qualquer sistema de informação hospitalar real;

    \item \textbf{Anonimato de usuários em pesquisas}: Quaisquer análises
    de desempenho pedagógico realizadas a partir dos dados do sistema
    devem respeitar o anonimato dos discentes, em conformidade com o
    Artigo 7.º, Inciso III da LGPD;

    \item \textbf{Consentimento Livre e Esclarecido}: A utilização do
    sistema em estudos de avaliação de aprendizagem com turmas reais
    exigirá, obrigatoriamente, Termo de Consentimento Livre e Esclarecido
    (TCLE) aprovado pelo Comitê de Ética em Pesquisa (CEP) da UFS
    (CAAE a ser registrado em trabalho futuro).
\end{enumerate}

% ===========================================================================
\chapter{Glossário de Termos Técnicos e Clínicos}
\label{anex_glossario}
% ===========================================================================

\begin{quadro}[htb]
\centering
\small
\caption{Glossário de termos técnicos e clínicos do SIGEA-GTT.}
\label{qua_glossario}
\begin{tabularx}{\textwidth}{l X}
\toprule
\textbf{Termo} & \textbf{Definição} \\
\midrule
\textbf{EA} & Evento Adverso: dano não intencional causado por intervenção
  de saúde e não pela condição subjacente do paciente \cite{who2009patient}. \\
\midrule
\textbf{GTT} & \textit{Global Trigger Tool}: ferramenta retrospectiva do IHI
  para identificação de eventos adversos em prontuários clínicos por meio de
  gatilhos clínicos padronizados \cite{griffin2009ihi}. \\
\midrule
\textbf{IHI} & \textit{Institute for Healthcare Improvement}: organização
  internacional sem fins lucrativos sediada em Boston (EUA), líder mundial
  em iniciativas de melhoria da segurança e qualidade em saúde. \\
\midrule
\textbf{JWT} & \textit{JSON Web Token}: padrão aberto (RFC 7519) para
  transmissão segura de informações entre partes como um objeto JSON
  assinado digitalmente \cite{jwt_rfc7519}. \\
\midrule
\textbf{LGPD} & Lei Geral de Proteção de Dados (Lei n.º 13.709/2018):
  lei brasileira que regulamenta o tratamento de dados pessoais em
  território nacional \cite{lgpd2018}. \\
\midrule
\textbf{NCC MERP} & \textit{National Coordinating Council for Medication
  Error Reporting and Prevention}: órgão norte-americano responsável pela
  taxonomia de categorias de gravidade de dano por eventos medicamentosos. \\
\midrule
\textbf{PDCA} & \textit{Plan-Do-Check-Act}: ciclo de melhoria contínua
  desenvolvido por Walter Shewhart e popularizado por W. Edwards Deming. \\
\midrule
\textbf{PNSP} & Programa Nacional de Segurança do Paciente: programa
  do Ministério da Saúde do Brasil instituído pela Portaria GM/MS
  n.º 529/2013, alinhado às Metas Internacionais de Segurança do Paciente
  da OMS \cite{anvisa2017}. \\
\midrule
\textbf{RBAC} & \textit{Role-Based Access Control}: modelo de controle de
  acesso baseado em funções/perfis de usuário, implementado no SIGEA-GTT
  com os perfis ADMIN, PROFESSOR e STUDENT. \\
\midrule
\textbf{SUS} & Sistema Único de Saúde: sistema público de saúde do Brasil,
  responsável pelo atendimento integral e universal à população, no qual
  o HU-UFS/EBSERH está inserido. \\
\midrule
\textbf{$TI_{1000}$} & Taxa de Incidência de Eventos Adversos por 1.000
  pacientes-dia: indicador epidemiológico do IHI-GTT calculado como
  $TI_{1000} = \frac{N_{EA}}{P_{dia}} \times 1000$. \\
\midrule
\textbf{$TI_{100}$} & Taxa de Incidência de Eventos Adversos por 100
  admissões: indicador IHI calculado como
  $TI_{100} = \frac{N_{EA}}{N_{adm}} \times 100$. \\
\bottomrule
\end{tabularx}
\fonte{Elaborado pelo autor (2026).}
\end{quadro}

\end{anexosenv}
